% Part: first-order-logic
% Chapter: sequent-calculus
% Section: derivations

\documentclass[../../../include/open-logic-section]{subfiles}

\begin{document}

\iftag{FOL}
      {\olfileid{fol}{seq}{der}}
      {\olfileid{pl}{seq}{der}}

\olsection{\usetoken{P}{derivation}}

\begin{explain}
Kita wis nerangake wujude sekuen wiwitan lan wis menehi aturan
inferensi. !!^{derivation}s ing kalkulus sekuen dibangkitake kanthi
induktif saka rong perkara iku: saben !!{derivation} dadi sekuen
wiwitan dhewe, utawa dumadi saka siji utawa rong !!{derivation}s kang
diterusake dening sawijining inferensi.
\end{explain}

\begin{defn}[$\Log{LK}$ !!{derivation}]
\emph{$\Log{LK}$-!!{derivation}} saka sekuen~$S$ yaiku wit winates
saka sekuen kang nyukupi syarat-syarat iki:
\begin{enumerate}
\item Sekuen-sekuen paling dhuwur ing wit iku sekuen wiwitan.
\item Sekuen paling ngisor ing wit iku~$S$.
\item Saben sekuen ing wit kejaba $S$ iku premis saka panrapan aturan
  inferensi kang bener, lan kesimpulane mapan langsung ing sangisore
  sekuen kasebut ing wit.
\end{enumerate}
Banjur kita ngarani $S$ \emph{sekuen-pungkasan} saka !!{derivation}
lan kandha yen $S$ \emph{bisa !!{derivable} ing $\Log{LK}$} (utawa
$\Log{LK}$-!!{derivable}).
\end{defn}

\begin{ex}
Saben sekuen wiwitan, upamane $!C \Sequent !C$, iku
!!a{derivation}. Kita bisa ngasilake !!{derivation} anyar saka iki
kanthi ngetrapake, umpamane, aturan $\LeftR{\Weakening}$,
\begin{prooftree}
\Axiom$ \Gamma \fCenter \Delta $
\RightLabel{\LeftR{\Weakening}}
\UnaryInf$ !A, \Gamma \fCenter \Delta$
\end{prooftree}
Nanging, aturan kasebut dimaksudake umum: $!A$ ing aturan bisa diganti
nganggo sembarang !!{sentence}, upamane uga nganggo~$!D$. Yen premis
cocog karo sekuen wiwitan $!C \Sequent !C$, iki ateges $\Gamma$ lan
$\Delta$ kalorone mung~$!C$, mula kesimpulane dadi $!D, !C \Sequent
!C$. Dadi, iki !!a{derivation}:
\begin{prooftree}
\Axiom$ !C \fCenter !C $
\RightLabel{\LeftR{\Weakening}}
\UnaryInf$ !D, !C \fCenter !C$
\end{prooftree}
Saiki kita bisa ngetrapake aturan liyane, upamane
$\LeftR{\Exchange}$, kang ngidini kita ngijolake rong !!{sentence}s ing
sisih kiwa. Dadi, iki uga !!{derivation} kang bener:
\begin{prooftree}
\Axiom$ !C \fCenter !C $
\RightLabel{\LeftR{\Weakening}}
\UnaryInf$ !D, !C \fCenter !C$
\RightLabel{\LeftR{\Exchange}}
\UnaryInf$ !C, !D \fCenter !C$
\end{prooftree}
Ing panrapan aturan iki, kang diwenehake kanthi wujud
\begin{prooftree}
\Axiom$ \Gamma, !A, !B, \Pi \fCenter \Delta $
\RightLabel{\LeftR{\Exchange}}
\UnaryInf$ \Gamma, !B, !A, \Pi \fCenter \Delta,$
\end{prooftree}
$\Gamma$ lan $\Pi$ kalorone kosong, $\Delta$ yaiku $!C$, lan peran
$!A$ lan $!B$ dilakoni dening $!D$ lan~$!C$, kanthi urut. Kanthi cara
kang meh padha, kita uga bisa weruh yen
\begin{prooftree}
\Axiom$ !D \fCenter !D $
\RightLabel{\LeftR{\Weakening}}
\UnaryInf$ !C, !D \fCenter !D$
\end{prooftree}
iku !!a{derivation}. Saiki rong !!{derivation}s iki bisa dijupuk lan
digabungake nganggo $\RightR{\land}$. Aturan kasebut yaiku
\begin{prooftree}
\Axiom$\Gamma \fCenter \Delta, !A$
\Axiom$ \Gamma \fCenter \Delta, !B$
\RightLabel{\RightR{\land}}
\BinaryInf$ \Gamma \fCenter \Delta, !A \land !B $
\end{prooftree}
Ing kasus iki, premis kudu cocog karo sekuen pungkasan saka
!!{derivation}s kang mungkasi premis. Iki ateges $\Gamma$ yaiku
$!C, !D$, $\Delta$ kosong, $!A$ yaiku $!C$, lan $!B$ yaiku $!D$. Mula,
yen inferensine bener, kesimpulane yaiku $!C, !D \Sequent !C
\land !D$.
\begin{prooftree}
\Axiom$ !C \fCenter !C $
\RightLabel{\LeftR{\Weakening}}
\UnaryInf$ !D, !C \fCenter !C$
\RightLabel{\LeftR{\Exchange}}
\UnaryInf$ !C, !D \fCenter !C$
\Axiom$ !D \fCenter !D $
\RightLabel{\LeftR{\Weakening}}
\UnaryInf$ !C, !D \fCenter !D$
\RightLabel{\RightR{\land}}
\BinaryInf$ !C, !D \fCenter !C \land !D $
\end{prooftree}
Mesthi wae, kita uga bisa ngowahi urutan premis, mula $!A$
dadi $!D$ lan $!B$ dadi~$!C$.
\begin{prooftree}
\Axiom$ !D \fCenter !D $
\RightLabel{\LeftR{\Weakening}}
\UnaryInf$ !C, !D \fCenter !D$
\Axiom$ !C \fCenter !C $
\RightLabel{\LeftR{\Weakening}}
\UnaryInf$ !D, !C \fCenter !C$
\RightLabel{\LeftR{\Exchange}}
\UnaryInf$ !C, !D \fCenter !C$
\RightLabel{\RightR{\land}}
\BinaryInf$ !C, !D \fCenter !D \land !C $
\end{prooftree}
\end{ex}

\end{document}
